% =============================================================
% Conclusion
% =============================================================

\section{Conclusion}
\label{sec:conclusion}

% 中文: 核心综合——brewing 揭示 transformer 深度的双重角色：计算答案 vs. 将答案重新格式化为可解码形式。这不是同一个过程，且后半段才是主要失败来源。
The central lesson of the brewing-to-resolution framework is that transformer depth serves two separable functions in code reasoning: \emph{computing} the answer and \emph{reformatting} it into a self-decodable state. These are not the same process---probing recovers the answer around 20\% depth while CSD catches up only near 72\%---and the gap between them is where most of the network's work happens. That the brewing scaffold remains stable across architectures and scales while resolution success changes with capability suggests these two functions are governed by different aspects of the model: the residual-stream geometry determines \emph{where} the transition occurs, while learned skill determines \emph{whether} it succeeds.

% 中文: 这改变了我们对 overthinking 和 early exit 的理解——问题不是模型"想太多"，而是后层主动破坏已有的正确计算。Overprocessed 是跨任务的主导失败模式。
This reframes the overthinking problem. Prior work treats overthinking as excess computation on easy inputs \citep{halawi2024overthinking}; our causal experiments show it is better understood as late-layer corruption of already-correct computation---and it is the dominant failure mode across all six task families, including control-flow tasks that are otherwise well-handled. Early-exit strategies that rely on output-distribution confidence may therefore be insufficient: the model's own decoding head is precisely what destroys the answer.

% 中文: 开放问题——为什么后层会破坏？brewing 是残差流架构的结构性后果还是训练动态的产物？
Two questions remain genuinely open. First, \emph{why} do late layers corrupt correct answers---is Overprocessing driven by representational interference from the unembedding matrix, by attention heads that over-attend to surface cues, or by something else? The causal tools introduced here can localize the phenomenon but do not yet explain its mechanism. Second, the stability of brewing across architectures raises a deeper question: is the availability-before-readiness ordering a necessary consequence of how residual streams accumulate information, or a contingent property of current training regimes that could be eliminated?
